\section{Exercise 2: Implementing the BPE tokenization algorithm}


\subsection*{Goal}

The goal of this lab was to implement the Byte Pair Encoding (BPE) algorithm from
scratch, using only the Python standard library, and then to compare the result with
the \texttt{sentencepiece} library. The corpus is the French Wikipedia page
\emph{Grèce antique}, about 225{,}058 characters and 6{,}495 distinct words.

\subsection*{What was implemented}

\textbf{Preparation.} The text is cut into words with the regular expression
\verb|\b[^\s]+\b|, and each word is counted with a \texttt{Counter}. Each word is then
split into single characters, which gives the starting vocabulary: 106 distinct
characters for this corpus.

\textbf{The three functions of the algorithm.}
\begin{itemize}
  \item \texttt{compute\_pair\_freqs(splits, word\_freqs)} goes through every word and
        counts each pair of successive tokens, weighted by the frequency of the word.
  \item \texttt{most\_frequent(pair\_freqs)} returns the pair with the highest count.
  \item \texttt{merge\_pair(a, b, splits)} rebuilds every word with the pair
        \texttt{(a,b)} replaced by the single token \texttt{ab}. It scans each word with
        an index and jumps forward by two positions after a merge, so that the second
        token of the pair is not examined again.
\end{itemize}

\textbf{The \texttt{BPE} class} repeats these three steps until the vocabulary reaches
the requested size. At each iteration it stores the merge in \texttt{merge\_rules}, adds
the new token to the vocabulary, and updates the splits. The \texttt{tokenize} method
applies the same rules, in the same order, to a new text. With a target size of 500, the
algorithm learns 394 merge rules on top of the 106 initial characters. The first merges
are the most frequent French patterns: \texttt{es}, \mk\texttt{d}, \mk\texttt{l},
\texttt{nt}, \texttt{qu}, \texttt{re}, \texttt{on}, \mk\texttt{de}, \texttt{ent}.

\subsection*{Detokenization}

The first version of the tokenizer could not rebuild the original text, because the
spaces were removed by the regular expression and the words were flattened into a single
list. Undoing the merges does not help: a merged token is exactly the concatenation of
its two parts, so the letters are already there; what is missing is the position of the
word boundaries.

The solution was to make the boundary part of the text itself. A marker \mk{} (the
character \texttt{U+2581}, the one used by \texttt{sentencepiece}) is added at the
beginning of every word, both during training and during tokenization. BPE then learns
it like any other character, and produces tokens such as \mk\texttt{en} or
\mk\texttt{Grèce}. Detokenization becomes a single line: join the tokens and replace the
marker by a space. The sentence \emph{cultures grecques développée en Grèce} is
recovered exactly.

\subsection*{Comparison with sentencepiece}

Both vocabularies have 500 tokens and share 456 of them, and their first 20 merges are
identical, which confirms that the implementation is correct. The remaining differences
come from the preprocessing, not from the algorithm:

\begin{itemize}
  \item \texttt{sentencepiece} reserves three special tokens
        (\texttt{<unk>}, \texttt{<s>}, \texttt{</s>});
  \item it keeps only the characters covering $99.95\%$ of the text (83 instead of 106),
        which leaves room for 414 merges instead of 394;
  \item our regular expression deletes the punctuation (2{,}666 commas, 1{,}510 full
        stops, 829 parentheses), so it cannot be learned nor restored;
  \item it never merges a letter with a punctuation character, so it learns
        \texttt{époque} alone where our version learns \mk\texttt{l'époque}.
\end{itemize}

To obtain a similar vocabulary, the punctuation should be kept as separate chunks (for
example with \verb!\w+|[^\w\s]+!), rare characters should be replaced by \texttt{<unk>},
and the three special tokens should be added to the vocabulary.

\subsection*{Results}

The final tokenizer uses 2.52 tokens per word on average (2.52 characters per token).
Tokenizing a word of the corpus gives exactly the split learned during training, every
token produced belongs to the vocabulary, and the detokenization returns the original
text for both known and unknown words.

\end{document}
